% paper.tex — Flagging soiled rooftop PV from public data alone.
% Build: see README.md. Figures must be generated before compiling.
\documentclass[10pt,twocolumn]{article}

\usepackage[letterpaper,margin=0.85in]{geometry}
\usepackage[T1]{fontenc}
\usepackage[utf8]{inputenc}
\usepackage{times}
\usepackage{graphicx}
\usepackage{booktabs}
\usepackage{amsmath}
\usepackage{amssymb}
\usepackage{siunitx}
\usepackage{url}
\usepackage[numbers,sort&compress]{natbib}
\usepackage{caption}
\usepackage{subcaption}
\usepackage{xcolor}
\usepackage{microtype}
\usepackage[colorlinks=true,allcolors=blue!55!black]{hyperref}

\sisetup{detect-all,group-separator={,}}
\captionsetup{font=small,labelfont=bf}
\graphicspath{{build/}}

\newcommand{\artifact}[1]{{\small\texttt{#1}}}

% A heading that wraps is justified by default, which stretches its interword
% space across the whole column. Run-in \paragraph titles are the worst affected
% because they are the longest. Ragged right on all three levels.
\makeatletter
\renewcommand\section{\@startsection{section}{1}{\z@}%
  {-3.5ex \@plus -1ex \@minus -.2ex}{2.3ex \@plus.2ex}%
  {\normalfont\Large\bfseries\raggedright}}
\renewcommand\subsection{\@startsection{subsection}{2}{\z@}%
  {-3.25ex\@plus -1ex \@minus -.2ex}{1.5ex \@plus .2ex}%
  {\normalfont\large\bfseries\raggedright}}
% The negative afterskip in article's default makes \paragraph a RUN-IN heading:
% the title joins the body paragraph and is justified with it, so a title that
% nearly fills a line gets its interword space stretched across the column and
% \raggedright in the style above cannot reach it. These titles are sentences,
% not labels, so give them their own line (positive afterskip) instead.
\renewcommand\paragraph{\@startsection{paragraph}{4}{\z@}%
  {2.4ex \@plus1ex \@minus.2ex}{0.5ex \@plus .1ex}%
  {\normalfont\normalsize\bfseries\raggedright}}
\makeatother

\title{\bf Permissively Licensed Rooftop Array Detection at \SI{21}{\centi\metre},
with an Executable Gate\\
\large and Per-Array Roof Geometry from Public Lidar}

\author{
  Cameron Gordon \quad Joshua Ramirez \quad Akshitha Nagaraj \quad Craig Fellers\\[2pt]
  {\normalsize Better Behavior Foundation}\\[2pt]
  {\small\texttt{betterbehaviorfoundation@gmail.com} \quad
   \url{https://betterbehaviorfoundation.com}}
}
\date{}

\begin{document}
\maketitle

\begin{abstract}
\noindent
\textbf{Draft, not a submission.} This is the detection and roof-geometry half of a
paper split because it was three papers in fifteen pages. The material below is moved
from that draft and has not been rewritten to stand alone: it still refers forward to a
soiling model that now lives in the other paper, and the introduction, related work and
discussion are not written. It is released in this state so that the claims the companion
paper cites can be checked, not as a finished document. Author list matches the companion
because it is the same work by the same people. See \texttt{SPLIT\_PLAN.md}.

The contribution here is a detection gate stated as an executable definition rather
than a number, a mask stage whose failure mode is a prompt-box problem rather than a
mask problem, an independent recall probe against county building permits the
detector never saw, and a projection trap in lidar-derived tilt that biases the angle
by 18.9\% on web-Mercator point clouds.
\end{abstract}

\section{Methods}
\subsection{Array detection}
\label{sec:detection}

\paragraph{Imagery.}
Santa Cruz County publishes 2025 aerial orthoimagery as a free map service. We tile a
pilot area of interest over the Live Oak and Pleasure Point side of the city into 249
tiles covering \SI{14.93}{\kilo\metre\squared} of \emph{actually imaged} ground; the
enclosing bounding box is \SI{23.71}{\kilo\metre\squared}, and the distinction matters for
any recall denominator. True ground GSD is \SI{0.204}{\metre}. Tiles are stored in
EPSG:3857, where the nominal pixel size reads \SI{0.256}{\metre}, inflated by
$1/\cos(\SI{37}{\degree})$. Each tile records its CRS and affine transform, so any
pixel-space detection projects back to geographic coordinates unambiguously. The county
service natively serves \SI{6.3}{\centi\metre} and we under-requested by $3.3\times$
linear, so the ceiling on this method sits above what we report.

\paragraph{Labels and splits.}
Arrays are annotated as polygons, because array extent feeds the area-to-dollar chain and
a box overstates it on non-rectangular roofs. The dataset comprises 976 chips
(680 train / 100 validation / 196 test) and 3{,}894 chip-level polygons, rebuilt from 244
of the 249 tiles. \textbf{Two ground-truth counts exist and are not interchangeable}:
chip-level (validation 385, test 636) double-counts arrays that the $\sim$80-pixel chip
overlap repeats and splits those the seams divide, whereas tile-level (validation 340,
test 585) does not. The gate scores tile-level; a tile-level $F_1$ quoted against 636 is
wrong by construction.

\paragraph{Model and inference path.}
RF-DETR at 728-pixel input, single class. Production inference decomposes each full tile
into a $2\times2$ grid of 640-pixel chips, runs one forward pass per chip, and merges
across seams by non-maximum suppression at IoU 0.55.

\paragraph{The gate, stated as an executable definition.}
Our shipping criterion is \emph{whatever \artifact{scripts/detect/eval\_tile\_f1.py}
prints}: tile-level detection $F_1$ at box-IoU $\geq 0.50$, micro-averaged over the split,
measured through the production path above. Three protocol choices define it.

\begin{enumerate}\itemsep2pt
\item \textbf{Confidence is tuned on validation and frozen; test is the reported number.}
Tuning confidence on test spends the held-out split. The cost of not doing so is
measured rather than assumed: test's own best $F_1$ is 0.8431 against 0.8260 at the
validation-frozen threshold (Figure~\ref{fig:sweep}).
\item \textbf{Matching is on boxes, not masks}, so the gate is invariant to the presence of
the mask stage. SAM2 changes area accuracy, not what was detected; conflate the two and
a regression cannot be traced to its stage.
\item \textbf{The 95\% interval bootstraps over tiles, not objects.} Arrays within a tile
are correlated, and resampling objects yields an interval that is too narrow.
\end{enumerate}

Passing requires all three of: $F_1 \geq 0.75$; $F_1$ 95\% CI lower bound $\geq 0.70$; and
recall $\geq 0.70$. Recall carries its own condition because a missed array is an
unrecoverable missed lead, while the permit join filters false positives downstream.

\paragraph{Mask stage and prompt-box sensitivity.}
SAM2 is in the product path not because the detector needs it, since the gate is
box-based, but because \emph{area is the product}. Prompted with the unmodified detector
box, SAM2 returns whole-roof masks at a high rate. We treat the prompt box as the design
variable and sweep it (\S\ref{sec:results-mask}).

\subsection{Independent recall probe}
\label{sec:permits}

Recall measured against held-out labels overstates recall against the real install base,
because labels and model see the same imagery. We therefore cross-check detections against
Santa Cruz County building-permit records, which the detector never trained on. Matching
is by parcel APN rather than geocoded point. The Census batch geocoder interpolates each
address onto the street segment, so points land a median \SI{22}{\metre} from the nearest
detected array and zero of 126 point matches fall inside a polygon. The point method
measures geocode quality, not detection.

\subsection{Roof geometry from public lidar}
\label{sec:geometry}

We fit a plane to USGS 3DEP~\cite{usgs3dep} lidar returns inside each detected polygon,
recovering per-array tilt and downslope azimuth over plain HTTP against the Entwine
octree. The headline 3DEP product is a bare-earth DEM and is useless here by construction,
since roofs are removed from it; the point cloud is not.

\paragraph{A projection trap.}
Entwine $X$/$Y$ are EPSG:3857 metres while $Z$ is true metres. At this latitude the
Mercator scale factor is $1/\cos(\SI{36.97}{\degree}) = 1.2515$, so a plane fitted on raw
coordinates reports tilt \textbf{18.9\% too shallow, silently, on every roof, producing an
entirely plausible distribution}. We de-inflate $X$/$Y$ by $\cos(\text{lat})$ and pin the
magnitude in a unit test, so the correction cannot be removed unnoticed. Anyone
reproducing this on web-Mercator lidar will hit it.

\paragraph{Failed fits are reported as null.} Arrays without a usable fit keep a flagged
flat default and are never median-filled---precisely the failure mode
\S\ref{sec:ranking} is about.


\section{Results}
\subsection{Detection}
\label{sec:results-detection}

Table~\ref{tab:gate} reports the gate. The shipping checkpoint passes all three
conditions. Figure~\ref{fig:sweep} shows the full operating curve on both splits and the
size of the tuning penalty.

\begin{table}[t]
\centering\small
\caption{Tile-level detection gate, box-IoU $\geq0.50$, micro-averaged over the test
split through the production chip-grid path. Confidence $c^*$ tuned on validation and
frozen. Intervals bootstrap over tiles ($n_{\text{tiles}}=49$,
$n_{\text{gt}}=585$).}
\label{tab:gate}
\begin{tabular}{lcc}
\toprule
 & \textbf{W2 (shipped)} & W1 (anchor) \\
\midrule
$c^*$ (val-frozen)      & 0.50 & 0.40 \\
Precision               & 0.8499 & 0.8533 \\
Recall                  & \textbf{0.8034} & 0.7556 \\
$F_1$                   & \textbf{0.8260} & 0.8015 \\
95\% CI                 & [0.798, 0.853] & [0.773, 0.830] \\
Test-tuned best $F_1$   & 0.8431 & 0.8044 \\
Gate verdict            & \textbf{PASS} & PASS \\
\bottomrule
\end{tabular}
\end{table}

\begin{figure*}[t]
  \centering
  \includegraphics[width=0.98\textwidth]{fig_threshold_sweep.pdf}
  \caption{Operating curves on both splits. The frozen threshold $c^*=0.50$ was chosen on
  validation alone; the test-tuned optimum sits 0.017 $F_1$ higher, which is the measured
  price of not spending the held-out split.}
  \label{fig:sweep}
\end{figure*}

\paragraph{Paired comparison, not overlapping intervals.}
Doubling the training set moved $F_1$ from 0.8015 to 0.8260. The two marginal intervals
overlap heavily and each point estimate lies inside the other's interval, so the naive
comparison reads ``within noise.'' That is the wrong test: both models were scored on the
same 49 tiles, so the comparison is paired. Bootstrapping the per-tile difference gives
$\Delta F_1 = +0.0246$, 95\% CI $[+0.0051, +0.0441]$, $P(\text{W2}>\text{W1}) = 0.993$,
concentrated in recall ($+0.0479$, CI $[+0.0202, +0.0753]$, $P=0.999$). The marginal
comparison would have discarded a real gain.

\paragraph{Two caveats on the shipped checkpoint.}
The tuning penalty grew relative to the anchor (0.017 against 0.003), so validation and
test disagree more about where the operating point sits; and the chip-to-tile gap widened
(0.029 against 0.008). Both are tracked and neither is yet explained.

\paragraph{Independent recall.}
Against 434 geocoded permit points inside the imaged footprint, APN-matched recall is
\textbf{55.3\%} (240/434) across all years and \textbf{73.6\%} (156/212) restricted to
pre-cutoff installs that the imagery could contain. The point method reads 29.0\%
(Wilson CI 25.0--33.5\%) for the geocoding reason given in \S\ref{sec:permits}. Three
caveats travel with the number. Part of the gain could in principle come from spurious
detections, since APN recall counts a parcel as found if any detected polygon's interior
point falls inside it; gate precision of 0.850 bounds but does not eliminate this.
Permits are a valid recall probe and an \emph{invalid} precision denominator: 85\% of
hand-labelled arrays sit on parcels with no permit on file, because the registry begins in
2016 and the residential boom preceded it. And 1{,}310 of 1{,}865 sites fall in city APN
books with 0.0\% county-permit coverage, because the City of Santa Cruz is a separate
permitting authority. Any study joining detections to permits must check jurisdiction
before reporting coverage.

\subsection{Mask quality is a prompt-box problem}
\label{sec:results-mask}

Figure~\ref{fig:sam} sweeps the prompt box over 60 validation arrays. Prompting with the
raw detector box yields median IoU 0.509, median mask area $1.97\times$ ground truth, and
a 43.3\% rate of whole-roof grabs. Shrinking the box 15\% yields median IoU \textbf{0.844},
median area ratio \textbf{1.01}, and \textbf{0.0\%} roof-grab. Area is therefore
essentially unbiased, which is exactly what the area-to-kW-to-dollar chain needs.

\begin{figure*}[t]
  \centering
  \includegraphics[width=0.98\textwidth]{fig_sam_promptbox.pdf}
  \caption{SAM2 prompt-box sensitivity. The whole-roof failure mode is created by an
  oversized prompt, not by the mask model, and disappears at the production default.}
  \label{fig:sam}
\end{figure*}

\paragraph{Two principled alternatives, both measured, both rejected.}
Background negative-point prompts behave exactly as theory predicts: they help when the
box is bad (roof-grab 38.3\%$\rightarrow$26.7\% at dilate-50\%) and hurt when it is good
(0.844$\rightarrow$0.820, introducing 1.7\% roof-grab where there was none), because with
an already-shrunk box the ring sits near the true panel edge and clips real pixels.
Multimask candidates clipped to the detector box with area-fit selection were worse
everywhere (0.824$\rightarrow$0.673 on exact boxes; roof-grab 1.7\%$\rightarrow$71.7\% at
dilate-25\%), because picking the candidate whose area matches the detector box amplifies
box error in exactly the regime the change targeted.

\paragraph{A ceiling we cannot locate.}
We never measured inter-annotator agreement, because our annotation platform deduplicates
and one tile cannot go to two labellers. If two humans agree to only $\sim$0.85 IoU
on a 26-pixel array, then 0.844 is the label-noise floor and further mask work is
unmeasurable. We do not know which side of that line we are on.

\subsection{Roof geometry}
\label{sec:results-geometry}

Plane fits succeed for 3{,}068 of 3{,}362 arrays (91.3\%), with median RMS residual
\SI{0.040}{\metre} over $\sim$270 points per array; synthetic planes are recovered to
$<\SI{0.2}{\degree}$ tilt and $<\SI{1}{\degree}$ azimuth. Figure~\ref{fig:tilt} compares
the recovered distribution against installer-reported tilt from 8{,}573 utility
interconnection filings for the same county---a different instrument, a different decade,
and no shared failure mode. \textbf{Both medians are \SI{19.0}{\degree}}. The low tails
differ for a documented reason: the filings encode ``not reported'' as \SI{0}{\degree} and
we drop those rows, so the filings understate the shallow tail (0.78\% below
\SI{5}{\degree}) while the lidar tail stays unbiased (8.38\%). That 8.38\% is directly
relevant to \citet{mejia2013}'s finding on sub-five-degree arrays.

\begin{figure*}[t]
  \centering
  \includegraphics[width=0.98\textwidth]{fig_tilt_validation.pdf}
  \caption{Per-array tilt from public lidar against installer-reported tilt from utility
  interconnection filings. The medians agree to \SI{0.03}{\degree}; the distributions do
  not, and the reason is reporting granularity rather than measurement error.}
  \label{fig:tilt}
\end{figure*}

\paragraph{What this comparison can and cannot establish.}
It is distributional, not paired, and that is a property of the source rather than a
choice: the DGStats interconnection extract carries a service city and nothing finer---no
parcel number, address or coordinate---so no lidar array can be matched to the filing that
describes it. 3D-PV-Locator~\cite{mayer2022pvlocator} validates per system and is the
stronger design where such a key exists.

We also decline to read agreement from the medians alone. A two-sample
Kolmogorov--Smirnov test rejects decisively ($D=0.214$, $p<10^{-90}$, $n=3{,}068$ against
$8{,}573$): the lidar distribution has much heavier tails, with a 10th percentile of
\SI{6.4}{\degree} against \SI{14}{\degree} reported and a 90th of \SI{30.8}{\degree}
against \SI{27}{\degree}. The mechanism is visible in the reported values. Every one is an
integer, five values account for 57.6\% of filings, and \textbf{27\% are exactly
\SI{18}{\degree}}---the 4:12 roof pitch. Installers file a nominal pitch from a small
menu; the lidar fit measures a continuous angle on whatever roof is actually there. That
the two medians land \SI{0.03}{\degree} apart across such different generating processes
is the real result, and it is a statement about central tendency only. Per-array evidence,
where we have it, comes instead from the standoff natural experiment below.

A natural experiment supports the interpretation: the lidar flight predates some
installations and postdates others, and permit dates split the population on a variable
unrelated to lidar. The panel standoff step appears exactly where physics says it should: pre-flight median
$+\SI{0.145}{\metre}$ against post-flight $+\SI{0.013}{\metre}$, difference
$+\SI{0.132}{\metre}$, bootstrap 95\% CI $[+0.084, +0.208]$, $p<0.0001$. The fit therefore
reports true racking angle where the array predates the flight and roof pitch where it
does not. The groups are small ($n=37$ and $54$) and only 15.9\% of arrays yield a usable
standoff.


\bibliographystyle{plainnat}
\bibliography{refs}
\end{document}
